\documentclass{article}
\usepackage{graphicx} % Required for inserting images
\usepackage{caption} % Required for \captionof
\usepackage{amsmath}
\usepackage{float}
\usepackage{comment}
\usepackage[T1]{fontenc}
\usepackage[utf8]{inputenc}
\usepackage{lmodern}
\usepackage{microtype}

\title{Hőtani mérés}
\author{Kern Luca, Csongor Vincze}
\date{2026 Április 1}

\begin{document}

\maketitle

\section*{Absztrakt}

Ebben a mérésben hőtani méréseket végeztünk, melynek során többek között meghatároztuk a víz fajhőjét elég nagy pontossággal, a használt rozsdamentes acél termoszok hőkapacitását, illetve más egyebeket. A mérőműszereink közt fontos szerepet játszott a vékony platina ellenállás-hőmérséklet összefüggésén alapuló digitális hőmérő.

\section*{1. Feladat: A hőmérsékleti adatok leolvasása digitális multiméterről}

Ebben a feladatban a multiméteres hőmérsékletleolvasással ismerkedtünk meg. A táblázatban szereplő együtthatók és a képlet alapján meghatároztuk, hogy a mért ellenállások mekkora hőmérsékletnek felelnének meg.

Alább találhatóak az adataink, amelyek azt mutatják be, hogy mekkorának mértük a szobahőmérsékletet a két különböző hőmérővel. Mint láthatjuk, ezek az adatok eltérnek. Ennek oka az, hogy az ellenállásos hőmérésnél egy eltolást okoz a kapcsolás soros ellenállása, ami hozzáadódik a hőmérő saját ellenállásához. Ez egy szisztematikus hiba. Ezt a következő feladatban tárgyaljuk részletesebben. Emellett természetesen a két Pt100 érzékelő sem teljesen azonos, és a kiértékeléshez használt képlet együtthatói is kismértékben eltérhetnek egymástól.
A mérést és kiértékelést a következőképpen folytattuk: Bedugtuk a multiméterbe a Pt100 szenzort, és az Easy DMM szoftverrel elindítottunk egy mérési sorozatot. Ezt kiexportáltuk egy csv fájlba. A csv fájlt beolvastuk, és a feladatban megadott képlet szerint (feketére: $ R[\Omega]=100 \Omega+0.385 \frac{\Omega}{{ }^{\circ} \mathrm{C}} \cdot T\left[{ }^{\circ} \mathrm{C}\right] $, zöldre (3-as termosz): $ R[\Omega]=100 \Omega+0.425 \frac{\Omega}{{ }^{\circ} \mathrm{C}} \cdot T\left[{ }^{\circ} \mathrm{C}\right] $) visszaszámoltuk a hőmérsékleteket, és ezt ábrázoltuk. Ezt mutatja az alábbi grafikon.

% zöld termosz:\\
% $R_{zold}=110.4 \Omega$\\
% $T_{zold}= \text{ 24.5°C}$\\
% fekete termosz:\\
% $R_{fekete}=111.3 \Omega$\\
% $T_{fekete}= \text{ 26.6°C}$\\


\begin{figure}[H]
    \centering
    \includegraphics[width=0.8\textwidth]{1_1_feladat.pdf}
    \caption*{1. feladat: Mért hőmérséklet az idő függvényében}
    \label{}
\end{figure}

A két platinahőmérő közel ugyanazt a hőmérsékletet mérte (27.8 és 27.9 °C), ami magasabb, mint a szobahőmérséklet, viszont ebbe még nincsen beleszámítva a hozzáadódó soros ellenállás. Amennyiben ezzel az értékkel kompenzálunk, a hőmérsékleti értékek csökkenni fognak.

Megjegyzés: Ebben a feladatban a labor alatt összekevertük a termoszok adatait, így a naplóban még más értékeket mutat a grafikon. Emellett itt levágtuk azt a részt, ahol kontakthiba (vagy más mérési zavaró tényező) miatt hibás értékeket mutatott a multiméter. Így csak az első 25 másodpercet ábrázoltuk, de ez bőven elég volt a hőmérők működésének validálására. Ahogy az ábra is mutatja, a fekete termosz ingadozása nagyon csekely, és a zöld termosz is megbízhatóan mért (az ábrázolt 25 s-os szakaszban kicsit emelkedett a zöld termosz szenzora által mért hőmérséklet, de ez igen kis mértékű volt. Utána valamilyen zavaró tényezőnek köszönhetően kicsit lecsökkent a zöld termosz hőmérséklete).

\section*{2. Feladat: A platinahőmérő kalibrálása}

Ebben a feladatban a hőmérők kalibrálását végeztük el, amire azért volt szükség, mert ahogy a hőmérők bekötésre kerültek, a saját ellenállásukhoz hozzáadódott egy soros ellenállás is a vezetékeknek és a banándugóknak köszönhetően. Ez a mért hőmérsékleti érték eltolódásához vezetett volna kalibrálás hiányában. A kalibrálást a következőképpen végeztük: mindkét termoszba hideg vizet töltöttünk úgy, hogy a víz ellepje a hőmérőt. Ezután a hőmérsékleti egyensúly beállta után megmértük a termoszokban a víz hőmérsékletét higanyos hőmérővel, és ugyanekkor megmértük a platinahőmérővel is. A mért hőmérsékletek különbségéből visszaszámoltuk a soros ellenállásokat.

Megjegyzés: A második alkalommal mi kaptuk meg a zöld termosz helyett a pirosat, így a második mérési alkalommal a piros termoszban lévő hőmérőhöz adódó ellenállást is meghatároztuk, ezért van három sor a táblázatunkban.

\begin{table}[H]
    \centering
    \renewcommand{\arraystretch}{1.2}
    \begin{tabular}{|l|c|c|c|}
        \hline
        Termosz & Hőmérővel mért hőmérséklet [°C] & Mért ellenállás [$\Omega$] & Soros ellenállás [$\Omega$] \\ \hline
        Zöld & 22,1 & 110,7 & 1,3 \\ \hline
        Fekete & 22,0 & 111,4 & 2,9 \\ \hline
        Piros & 22,4 & 110,29 & 1,7 \\ \hline
    \end{tabular}
    \caption{Termoszok hőmérsékleti korrekciója}
    \label{tab:hotan_2}
\end{table}

\begin{comment}
Termosz színkódja: zöld (3)\\
Egyensúlyi hőmérséklet: 22,1 °C\\
Mért ellenállás: 110,7 $\Omega$\\
Soros ellenállás: 1,3 $\Omega$\\

Termosz színkódja: fekete\\
Egyensúlyi hőmérséklet: 22,0 °C\\
Mért ellenállás: 111,4 $\Omega$\\
Soros ellenállás: 2,9 $\Omega$\\

Termosz színkódja: piros\\
Egyensúlyi hőmérséklet: 22,4 °C\\
Mért ellenállás: 110,29 $\Omega$\\
Soros ellenállás: 1,7 $\Omega$\\
\end{comment}

Az eredményeinkből látszik, hogy az egyes termoszokban lévő hőmérők valóban kicsit eltérően mérnek. Emellett az is megfigyelhető, hogy a soros ellenállások egymástól csak kismértékben térnek el.


\section*{3. Feladat: A termosz hőveszteségének mérése}

Ebben a feladatban azt mértük meg, hogy ha egy termoszt megtöltünk forró vízzel akkor mekkora a hőveszteség. Ennek a hőveszteségnek két fő oka lehet. Az egyik az, hogy maga a termosz is átmelegszik (ezt a következő feladatban tárgyaljuk részletesebben), illetve hogy a környezetnek is lead a rendszert valamennyi hőt. Mivel a víz és a termosz közti termikus egyensúly beállta után kezdtünk el mérni, így ezzel a méréssel azt tudtukmegállapítani, hogy a rendszert mennyi hőt ad át a környezetének. 5 percig gyűjtöttünk hőmérséklet adatokat, majd ábrázoltuk ezeket az adatokat az eltelt idő függvényében. Erre a grafikonra egyenest illesztettünk, és meghatároztuk a meredekségét, amiből meghatároztuk,hogy egy perc alatt nagyjából mekkora a hőveszteség. Ez $\Delta T=1$ °C-nak adódott.

\begin{figure}[H]
    \centering
    \includegraphics[width=0.8\textwidth]{3_1_graf.pdf}
    \caption*{3. feladat: A termosz hővesztesége az idő függvényében}
    \label{fig:3_1}
\end{figure}

A grafikonról leolvasható, hogy kicsi a hőveszteség, a meredekség szinte 0, ami megfelel a várakozásainknak.
Hőmérsékleti korrekció: .........        %kell még!


\end{document}